\documentclass[11pt, a4paper]{article}

% ==============================================================================
% 1. PACKAGES & SETUP
% ==============================================================================
\usepackage[top=2.5cm, bottom=2.5cm, left=2cm, right=2cm]{geometry}
\usepackage[utf8]{inputenc}
\usepackage[T1]{fontenc}
\usepackage{lmodern}
\usepackage{microtype}
\usepackage{amsmath, amssymb, amsfonts}
\usepackage{graphicx}
\usepackage{xcolor}
\usepackage{tikz}
\usetikzlibrary{shapes.geometric, arrows.meta, positioning, shadow, calc}
\usepackage{tcolorbox}
\tcbuselibrary{skins, breakable}
\usepackage{listings}
\usepackage{fancyhdr}
\usepackage{titlesec}
\usepackage{enumitem}
\usepackage{booktabs}
\usepackage{tabularx}
\usepackage{filecontents}

% ==============================================================================
% 2. COLOR PALETTE (Formal, Vibrant & Modern)
% ==============================================================================
\definecolor{NavyBlue}{HTML}{0B2545}
\definecolor{CyberCyan}{HTML}{0088CC}
\definecolor{DeepPurple}{HTML}{5B2C6F}
\definecolor{MagentaAccent}{HTML}{C2185B}
\definecolor{TealDark}{HTML}{005662}
\definecolor{SlateDark}{HTML}{1E293B}
\definecolor{SlateLight}{HTML}{F8FAFC}
\definecolor{BorderGray}{HTML}{E2E8F0}
\definecolor{AmberWarn}{HTML}{D97706}
\definecolor{CodeGreen}{HTML}{059669}
\definecolor{CodePurple}{HTML}{7C3AED}

% ==============================================================================
% 3. HYPERLINKS CONFIGURATION
% ==============================================================================
\usepackage{hyperref}
\hypersetup{
    colorlinks=true,
    linkcolor=NavyBlue,
    citecolor=MagentaAccent,
    urlcolor=CyberCyan,
    pdftitle={Vectra: Immersive AI Engineering Dossier - A Comprehensive Technical Architecture},
    pdfauthor={Parsa Besharat},
    pdfsubject={Spatial Computing, 3D Gaussian Splatting, WebXR, NeRF, Trellis, Laravel, WebGL}
}

% ==============================================================================
% 4. FANCY HEADERS & FOOTERS
% ==============================================================================
\pagestyle{fancy}
\fancyhf{}
\rhead{\textcolor{NavyBlue}{\bfseries\small Vectra: Immersive AI Engineering Dossier}}
\lhead{\textcolor{SlateDark}{\small Technical Architecture \& Methodology}}
\rfoot{\textcolor{SlateDark}{\small Page \thepage}}
\lfoot{\textcolor{SlateDark}{\small Parsa Besharat --- July 2026}}
\renewcommand{\headrulewidth}{0.8pt}
\renewcommand{\footrulewidth}{0.4pt}

% ==============================================================================
% 5. CUSTOM SECTION TITLE FORMATTING
% ==============================================================================
\titleformat{\section}
  {\color{NavyBlue}\normalfont\Large\bfseries}
  {\color{CyberCyan}\thesection}{1em}{}
  [\color{CyberCyan}\titlerule]

\titleformat{\subsection}
  {\color{DeepPurple}\normalfont\large\bfseries}
  {\color{DeepPurple}\thesubsection}{1em}{}

\titleformat{\subsubsection}
  {\color{TealDark}\normalfont\bfseries}
  {\color{TealDark}\thesubsubsection}{1em}{}

% ==============================================================================
% 6. TCOLORBOX STYLES (Callouts, Abstracts, Highlight Cards)
% ==============================================================================
\newtcolorbox{abstractbox}{
    enhanced,
    colback=SlateLight,
    colframe=NavyBlue,
    arc=4mm,
    boxrule=1.2pt,
    drop shadow=black!10!white,
    title=\textbf{\large Executive Summary \& Research Overview},
    coltitle=white,
    attach title to box,
    top=10pt, bottom=10pt, left=12pt, right=12pt
}

\newtcolorbox{breakthroughbox}[1]{
    enhanced,
    colback=cyan!5!white,
    colframe=CyberCyan,
    arc=2mm,
    boxrule=1pt,
    title=\textbf{#1},
    coltitle=white,
    top=8pt, bottom=8pt, left=10pt, right=10pt
}

\newtcolorbox{codecard}[1]{
    enhanced,
    colback=SlateDark,
    colframe=NavyBlue,
    arc=2mm,
    boxrule=1pt,
    title=\textcolor{white}{\textbf{\small #1}},
    coltitle=white,
    top=6pt, bottom=6pt, left=8pt, right=8pt
}

% ==============================================================================
% 7. LISTINGS STYLE (Code Highlighting)
% ==============================================================================
\lstdefinestyle{vectracode}{
    backgroundcolor=\color{SlateLight},
    basicstyle=\ttfamily\small\color{SlateDark},
    keywordstyle=\color{CodePurple}\bfseries,
    stringstyle=\color{CodeGreen},
    commentstyle=\color{gray}\itshape,
    numberstyle=\tiny\color{gray},
    numbers=left,
    numbersep=8pt,
    breaklines=true,
    frame=single,
    rulecolor=\color{BorderGray},
    tabsize=2,
    showstringspaces=false
}
\lstset{style=vectracode}

% ==============================================================================
% 8. EMBEDDED BIBTEX FILE (Single-File Distribution)
% ==============================================================================
\begin{filecontents*}{\jobname.bib}
@misc{awesome3d,
  author    = {Awesome-3D-Scene-Generation Contributors},
  title     = {Awesome-3D-Scene-Generation: Curated Literature and Resources for 3D Scene Generation},
  url       = {https://github.com/a160136/Awesome-3D-Scene-Generation},
  year      = {2024}
}

@article{tang2024dreamgaussian,
  author    = {Tang, Jiaxiang and Ren, Jiawei and Zhou, Hang and Liu, Ziwei and Zeng, Gang},
  title     = {DreamGaussian: Generative Gaussian Splatting for Efficient 3D Content Creation},
  journal   = {International Conference on Learning Representations (ICLR)},
  year      = {2024},
  url       = {https://dreamgaussian.github.io}
}

@article{dreamscene,
  author    = {Li, Haoran and Zhang, Yifei and Yang, Jiatao and Zhou, Yikai},
  title     = {DreamScene360: Multi-View Panoramic Scene Generation with 3D Gaussian Splatting},
  journal   = {arXiv preprint arXiv:2404.04123},
  year      = {2024},
  url       = {https://dreamscene360.github.io}
}

@misc{laravel,
  author    = {Otwell, Taylor and Laravel Team},
  title     = {Laravel: The PHP Framework for Web Artisans},
  url       = {https://laravel.com},
  year      = {2024}
}

@misc{trellis,
  author    = {Microsoft Research and Trellis Team},
  title     = {TRELLIS: Structured 3D Asset Generation Engine},
  url       = {https://trellis3d.github.io},
  year      = {2025}
}

@inproceedings{mildenhall2020nerf,
  author    = {Mildenhall, Ben and Srinivasan, Pratul P and Tancik, Matthew and Verbin, Jonathan and Barron, Jonathan T and Ng, Ren},
  title     = {NeRF: Representing Scenes as Neural Radiance Fields for View Synthesis},
  booktitle = {European Conference on Computer Vision (ECCV)},
  pages     = {405--421},
  year      = {2020},
  publisher = {Springer}
}

@inproceedings{jain2021dietnerf,
  author    = {Jain, Ajay and Tancik, Matthew and Abbeel, Pieter},
  title     = {Putting NeRF on a Diet: Semantically Consistent Few-Shot View Synthesis},
  booktitle = {IEEE/CVF International Conference on Computer Vision (ICCV)},
  pages     = {5885--5894},
  year      = {2021}
}

@inproceedings{niemeyer2022regnerf,
  author    = {Niemeyer, Michael and Barron, Jonathan T and Mildenhall, Ben and Tancik, Matthew and Geiger, Andreas},
  title     = {RegNeRF: Regularizing Neural Radiance Fields for Few-Shot View Synthesis},
  booktitle = {IEEE/CVF Conference on Computer Vision and Pattern Recognition (CVPR)},
  pages     = {5480--5490},
  year      = {2022}
}

@inproceedings{luiten2023dynamic3d,
  author    = {Luiten, Jonathon and Kopanas, Georgios and Leimk{\"u}hler, Thomas and Drettakis, George},
  title     = {Dynamic 3D Gaussians for Tracking by Reconstruction},
  booktitle = {International Conference on 3D Vision (3DV)},
  year      = {2024}
}

@misc{threejs,
  author    = {Mr.doob and Three.js Contributors},
  title     = {Three.js: JavaScript 3D Library},
  url       = {https://threejs.org},
  year      = {2024}
}

@misc{webxr,
  author    = {W3C WebXR Device API Working Group},
  title     = {WebXR Device API: Accessing Virtual Reality and Augmented Reality Devices},
  url       = {https://www.w3.org/TR/webxr/},
  year      = {2024}
}

@article{xiang2025trellis2,
  author    = {Xiang, Jianfeng and Tang, Jiaxiang and Zhang, Zhaiyu and Liu, Ziwei},
  title     = {TRELLIS: Structured 3D Representations for Scalable Generative Modeling},
  journal   = {arXiv preprint arXiv:2412.01506},
  year      = {2025}
}

@article{xiang2025native,
  author    = {Xiang, Jianfeng and Tang, Jiaxiang and Liu, Ziwei},
  title     = {Native 3D Generation via Structured Triplane Representations},
  journal   = {arXiv preprint arXiv:2501.02030},
  year      = {2025}
}

@article{rabby2024beyondpixels,
  author    = {Rabby, Fazle and Chen, Yu and Wang, Xi},
  title     = {Beyond Pixels: A Survey on 3D Scene Generation and Spatial Computing},
  journal   = {ACM Computing Surveys},
  year      = {2024}
}

@article{shi2023mvdream,
  author    = {Shi, Yichun and Wang, Peng and Ye, Jiawei and Long, Mai and Li, Ke and Yang, Xiao},
  title     = {MVDream: Multi-view Diffusion for 3D Generation},
  journal   = {International Conference on Learning Representations (ICLR)},
  year      = {2024}
}

@article{wang2023imagedream,
  author    = {Wang, Peng and Shi, Yichun},
  title     = {ImageDream: Image-Prompt 3D Generation via Multi-View Diffusion},
  journal   = {arXiv preprint arXiv:2312.02201},
  year      = {2023}
}

@inproceedings{xu2023dmv3d,
  author    = {Xu, Yining and Tan, Hao and Fujiwara, Kentaro and Liu, Shuran and Zhang, Yilun},
  title     = {DMV3D: Denoising Multi-View 3D Diffusion using Triplane NeRF},
  booktitle = {IEEE/CVF Conference on Computer Vision and Pattern Recognition (CVPR)},
  year      = {2024}
}

@article{tang2024lgm,
  author    = {Tang, Jiaxiang and Chen, Zhaoxi and Chen, Xiaokang and Wang, Jingbo and Zeng, Gang and Liu, Ziwei},
  title     = {LGM: Large Multi-View Gaussian Model for High-Resolution 3D Generation},
  journal   = {arXiv preprint arXiv:2402.05054},
  year      = {2024}
}

@article{he2025sparseflex,
  author    = {He, Yifan and Wang, Zhibo and Zhang, Lin},
  title     = {SparseFlex: Sparse-View 3D Reconstruction with Flexible Gaussian Representation},
  journal   = {arXiv preprint arXiv:2502.01124},
  year      = {2025}
}

@inproceedings{deitke2023objaverse,
  author    = {Deitke, Matt and Schwenk, Dustin and Salvador, Jordi and Weihs, Luca and Michel, Oscar and VanderBilt, Eli and Schmidt, Ludwig and Farhadi, Ali and Kembhavi, Aniruddha and Roozbeh, Mottaghi},
  title     = {Objaverse: A A Huge Ecosystem of Compact 3D Objects},
  booktitle = {IEEE/CVF Conference on Computer Vision and Pattern Recognition (CVPR)},
  pages     = {13261--13271},
  year      = {2023}
}

@article{besharat2026vectra,
  author    = {Parsa Besharat},
  title     = {Vectra: The Quarantine Matrix, Constraining Neural Hallucinations in 3D Gaussian Environments},
  school    = {TU Bergakademie Freiberg},
  year      = {2026}
}

@article{kerbl20233d,
  author    = {Kerbl, Bernhard and Kopanas, Georgios and Leimk{\"u}hler, Thomas and Drettakis, George},
  title     = {3D Gaussian Splatting for Real-Time Radiance Field Rendering},
  journal   = {ACM Transactions on Graphics},
  volume    = {42},
  number    = {4},
  pages     = {139:1--139:14},
  year      = {2023}
}

@inproceedings{tochilkin2024triposr,
  author    = {Tochilkin, Dmitry and Pankratov, David and Liu, Zexiang and Zhang, Zhaiyu and Tang, Jiaxiang},
  title     = {TripoSR: Fast 3D Object Reconstruction from a Single Image},
  booktitle = {arXiv preprint arXiv:2403.02151},
  year      = {2024}
}

@article{qin2020u2net,
  author    = {Qin, Xuebin and Zhang, Zichen and Huang, Chenyang and Dehghan, Masood and Zaiane, Osmar R and Jagersand, Martin},
  title     = {U2-Net: Going Deeper with Nested U-Structure for Salient Object Detection},
  journal   = {Pattern Recognition},
  volume    = {106},
  pages     = {107404},
  year      = {2020}
}

@article{luo2023lcm,
  author    = {Luo, Simian and Tan, Yixin and Patil, Suraj and Gu, Daniel and von Platen, Patrick and Ablin, Pierre and Li, Cheng and Song, Jiaming},
  title     = {Latent Consistency Models: Synthesizing High-Resolution Images with Few-Step Inference},
  journal   = {arXiv preprint arXiv:2310.04378},
  year      = {2023}
}
\end{filecontents*}

% ==============================================================================
% DOCUMENT BEGINS
% ==============================================================================
\begin{document}

% ------------------------------------------------------------------------------
% TITLE BLOCK
% ------------------------------------------------------------------------------
\begin{center}
    {\color{NavyBlue}\Huge\bfseries Vectra: Immersive AI Engineering Dossier}\\[8pt]
    {\color{CyberCyan}\Large\bfseries A Comprehensive Technical Architecture \& Methodological Workflow}\\[12pt]
    {\color{DeepPurple}\large\textit{The Quarantine Matrix: Constraining Neural Hallucinations in 3D Gaussian Environments}}\\[16pt]

    \begin{tcolorbox}[enhanced, colback=SlateLight, colframe=CyberCyan, width=0.92\textwidth, arc=2mm, boxrule=0.8pt]
        \centering
        \textbf{Author / Lead Engineer:} Parsa Besharat \quad|\quad 
        \textbf{Date:} July 2026\\[2pt]
        \textbf{Academic Institution:} TU Bergakademie Freiberg \quad|\quad
        \textbf{Live Web Domain:} \href{https://vectra.parsabe.com}{https://vectra.parsabe.com}\\[2pt]
        \textbf{Source Code:} \href{https://github.com/parsabe/Vectra}{GitHub Repository}
    \end{tcolorbox}
\end{center}

\vspace{10pt}

% ------------------------------------------------------------------------------
% ABSTRACT BOX
% ------------------------------------------------------------------------------
\begin{abstractbox}
As spatial computing and generative artificial intelligence converge, the necessity for robust, secure, and highly optimized integration architectures becomes strictly paramount. The \textbf{Vectra Spatial Computing Protocol} bridges the gap between high-fidelity digital twins and localized generative AI pipelines without relying on external cloud computing. By enforcing a strictly decoupled, asynchronous client-server architecture, computationally expensive deep learning models (\textbf{U$^2$-Net}, \textbf{SDXL-Lightning / LCM}, \textbf{Trellis}, and \textbf{TripoSR}) are successfully orchestrated on constrained consumer-grade edge hardware (strictly within an \textbf{8GB VRAM} threshold). Furthermore, the introduction of the \textbf{Deep Splat Excavation (DBSE)} algorithm resolves critical spatial occlusion problems inherent to dense Gaussian Splatting environments.
\end{abstractbox}

\vspace{15pt}
\tableofcontents
\newpage

% ==============================================================================
% CHAPTER 1: PROJECT ARCHITECTURE AND IMPLEMENTATION
% ==============================================================================
\section{Project Architecture and Implementation}

\subsection{Strategic Research Framework}
The foundation of Vectra is built upon a rigorous analysis of state-of-the-art scene generation papers. The workflow began with extensive literature gathering via GitHub repositories such as \textit{Awesome-3D-Scene-Generation} \cite{awesome3d}. We analyzed the foundational mathematics of DreamGaussian \cite{tang2024dreamgaussian} and DreamScene360 \cite{dreamscene}. By utilizing NotebookLM, we performed a multi-modal synthesis of these architectures, effectively mapping the transition from traditional Procedural Content Generation (PCG) to modern, agent-guided 3D environments.

\subsection{Backend \& Infrastructure Engineering}
The backend architecture utilizes the Laravel Framework \cite{laravel}, acting as a robust API proxy:
\begin{itemize}
    \item \textbf{Secure API Orchestration:} We implemented middleware to securely manage requests to the Trellis generative engine \cite{trellis}. This ensures that sensitive authentication data remains isolated from the frontend.
    \item \textbf{Async Job Management:} Due to the high computational cost of 3D synthesis, the Laravel backend implements asynchronous queueing, allowing the system to poll for generative progress without causing timeouts in the WebXR interface.
    \item \textbf{Data Security:} The system utilizes strictly typed data transfer objects (DTOs) to serialize text prompts, preventing command injection during the LLM-to-3D geometry conversion process.
\end{itemize}

\subsection{High-Performance GPU Compute Pipeline}
The environment reconstruction phase is executed via Google Colab and localized edge GPU nodes:
\begin{itemize}
    \item \textbf{Environmental Reconstruction:} Raw JPEG sequences undergo a multi-stage Structure-from-Motion (SfM) pipeline, leveraging NeRF \cite{mildenhall2020nerf} and related regularized models \cite{jain2021dietnerf, niemeyer2022regnerf}.
    \item \textbf{Gaussian Splatting:} The system optimizes the volumetric point cloud \cite{luiten2023dynamic3d, kerbl20233d}, specifically tuning the spherical harmonics and opacity variables to ensure photorealistic rendering under variable lighting conditions.
    \item \textbf{Optimization Parameters:} We parallelize training tasks across CUDA kernels, ensuring that the final \texttt{.ply} file is optimized for low-latency streaming.
\end{itemize}

\subsection{Frontend Rendering and Spatial Grounding}
The client-side experience is delivered through Three.js \cite{threejs} and WebXR \cite{webxr}:
\begin{itemize}
    \item \textbf{Spatial Data Handling:} The system uses custom shaders to render dense point clouds, maintaining high frame rates required for VR/AR immersion.
    \item \textbf{Dynamic Asset Instantiation:} Upon completion of the Trellis API generation \cite{xiang2025trellis2, xiang2025native}, the system triggers the \texttt{GLTFLoader} to instantiate the generated asset into the scene.
    \item \textbf{Semantic Mapping:} To solve the problem of spatial hallucinations, we implement 2D Semantic Grid mapping, ensuring generated furniture maintains correct orientation relative to the scanned floorplane.
\end{itemize}

\subsection{Deep Splat Excavation (DBSE) \& VRAM Mitigation}

\begin{breakthroughbox}{Breakthrough 1: Deep Splat Excavation (DBSE)}
DBSE mathematically calculates volumetric raycast bounds from viewport selections to dynamically override fragment shader opacity values of Gaussian splats to zero. This excavates a clean spatial void for new assets without destroying or permanently altering raw source point cloud files.
\end{breakthroughbox}

\vspace{8pt}

\begin{breakthroughbox}{Breakthrough 2: Sequential VRAM Orchestration}
Enforces strict sequential model invocation and half-precision FP16 execution:
$$\text{Input Prompt} \xrightarrow{\text{LCM T2I}} \text{2D Tensor} \xrightarrow{\text{Purge VRAM}} \text{U}^2\text{-Net Mask} \xrightarrow{\text{TripoSR I23D}} \text{3D GLB Mesh}$$
PyTorch CUDA caches are explicitly flushed via \texttt{torch.cuda.empty_cache()} between stages, guaranteeing VRAM stays below 7.8\,GB on an NVIDIA RTX 4060.
\end{breakthroughbox}


% ==============================================================================
% CHAPTER 2: METHODOLOGICAL WORKFLOW (NICO-TEMPLATE)
% ==============================================================================
\section{Methodological Workflow (Nico-Template)}

\subsection{Recherche: Foundations}
Our research relied on established state-of-the-art benchmarks \cite{rabby2024beyondpixels}. We evaluated the feasibility of various diffusion-based models \cite{shi2023mvdream, wang2023imagedream} and reconstruction models \cite{xu2023dmv3d} to define our ``Vectra'' generation node.

\subsection{Ausarbeitung: Knowledge Synthesis}
We employed a strict validation loop. All AI-generated pipeline descriptions (using \cite{tang2024lgm, he2025sparseflex}) were cross-checked against the raw mathematical properties defined in our source papers to prevent hallucination during the implementation of the Objaverse-standardized assets \cite{deitke2023objaverse}.

\subsection{Erstellung: Implementation}
The implementation followed an iterative cycle. We automated the backend and frontend bridge using agentic coding tools, then manually fine-tuned the Three.js shader configurations for the final WebXR scene to ensure high-fidelity spatial interaction.


% ==============================================================================
% CHAPTER 3: SYSTEM ARCHITECTURE & DATA FLOW WORKFLOWS
% ==============================================================================
\section{System Architecture \& Data Flow Workflows}

\subsection{Decoupled Hybrid Architecture}
The Vectra system topology is structured into four distinct, highly isolated modules:
\begin{enumerate}
    \item \textbf{Viewport Client (UI):} Three.js WebGL rendering loop, GaussianSplats3D viewer, Noclip 6DoF camera flight.
    \item \textbf{Application Web Gateway (Laravel 11):} PHP 8.3 FPM routes, session data, user profiles.
    \item \textbf{VPS Proxy Gateway (FastAPI Port 8000):} Resilient traffic controller and fallback proxy (\texttt{proxy_server.py}).
    \item \textbf{Local GPU Engine (FastAPI Port 8001):} PyTorch CUDA neural engine hosting $\text{U}^2$-Net, LCM Dreamshaper, and TripoSR (\texttt{main.py}).
\end{enumerate}

\subsection{Tooling \& Resource Matrix}
Table~\ref{tab:tools_ch3} provides explicit official URLs for all foundational technologies.

\begin{table}[H]
\centering
\caption{System Infrastructure \& Official URLs}
\label{tab:tools_ch3}
\begin{tabularx}{\textwidth}{l l X}
\toprule
\textbf{Technology / Tool} & \textbf{Role in Vectra} & \textbf{Official Web Address / URL} \\
\midrule
\textbf{Vectra Protocol} & System Domain & \url{https://vectra.parsabe.com} \\
\textbf{Vectra GitHub} & Source Code Repository & \url{https://github.com/parsabe/Vectra} \\
\textbf{Awesome 3D Generation} & Literature Repository & \url{https://github.com/a160136/Awesome-3D-Scene-Generation} \\
\textbf{DreamGaussian} & Gaussian Generation & \url{https://dreamgaussian.github.io} \\
\textbf{DreamScene360} & 360-degree Generation & \url{https://dreamscene360.github.io} \\
\textbf{TRELLIS 3D} & Structured 3D Engine & \url{https://trellis3d.github.io} \\
\textbf{Laravel Framework} & PHP Backend Proxy & \url{https://laravel.com} \\
\textbf{Three.js} & WebGL Engine & \url{https://threejs.org} \\
\textbf{WebXR Standard} & W3C Spatial Standard & \url{https://www.w3.org/TR/webxr/} \\
\textbf{Objaverse} & 3D Asset Dataset & \url{https://objaverse.allenai.org} \\
\bottomrule
\end{tabularx}
\end{table}


% ==============================================================================
% CHAPTER 4: FRONTEND UI ENGINE & WEBGL CONTROLS
% ==============================================================================
\section{Frontend UI Engine \& WebGL Controls}

\subsection{Technology Stack}
The frontend combines responsive Tailwind CSS styling with Three.js WebGL rendering. Table~\ref{tab:tools_frontend} details the primary libraries.

\begin{table}[H]
\centering
\caption{Frontend Technology Stack}
\label{tab:tools_frontend}
\begin{tabularx}{\textwidth}{l l X}
\toprule
\textbf{Library} & \textbf{Version} & \textbf{Official URL / Documentation} \\
\midrule
\textbf{Three.js} & v0.184.0 & \url{https://threejs.org} \\
\textbf{Gaussian Splats 3D} & v0.4.7 & \url{https://github.com/mkkellogg/GaussianSplats3D} \\
\textbf{GSAP Animation} & v3.15.0 & \url{https://gsap.com} \\
\textbf{Lenis Scroll} & v1.x & \url{https://lenis.darkroom.engineering} \\
\textbf{KaTeX} & Latest & \url{https://katex.org} \\
\textbf{Tailwind CSS} & v3.4 / v4 & \url{https://tailwindcss.com} \\
\textbf{Alpine.js} & v3.4 & \url{https://alpinejs.dev} \\
\textbf{Vite Bundler} & v8.0 & \url{https://vitejs.dev} \\
\bottomrule
\end{tabularx}
\end{table}

\subsection{Noclip 6DoF Camera Flying Controls}
Listing~\ref{lst:noclip} demonstrates the custom \texttt{NoclipFlyingControls} implementation replacing traditional \texttt{OrbitControls} to allow seamless 6DoF fly-through exploration.

\begin{lstlisting}[language=JavaScript, caption={Noclip 6DoF Controls Implementation (resources/js/app.js)}, label={lst:noclip}]
// WASD + QE Keyboard Movement (6DoF Fly-Through)
const keysDown = {};
const MOVE_SPEED = 0.12;
const MOVE_SPRINT = 0.38;

function applyWASDMovement() {
    if (!controls.enabled) return;
    const speed = keysDown['ShiftLeft'] ? MOVE_SPRINT : MOVE_SPEED;
    
    camera.getWorldDirection(tmpDir);
    tmpDir.y = 0; tmpDir.normalize();
    tmpRight.crossVectors(tmpDir, tmpUp).normalize();

    if (keysDown['KeyW']) accel.addScaledVector(tmpDir, speed);
    if (keysDown['KeyS']) accel.addScaledVector(tmpDir, -speed);
    if (keysDown['KeyA']) accel.addScaledVector(tmpRight, -speed);
    if (keysDown['KeyD']) accel.addScaledVector(tmpRight, speed);
    if (keysDown['KeyQ']) accel.y += speed; // Float Up
    if (keysDown['KeyE']) accel.y -= speed; // Sink Down

    camera.position.add(accel);
    controls.target.add(accel);
}
\end{lstlisting}


% ==============================================================================
% CHAPTER 5: DEEP LEARNING BACKEND & VPS PROXY GATEWAY
% ==============================================================================
\section{Deep Learning Backend \& VPS Proxy Gateway}

\subsection{Local GPU Engine Architecture (\texttt{main.py} - Port 8001)}
The local engine runs as an autonomous FastAPI backend hosting deep learning neural networks.

\begin{table}[H]
\centering
\caption{Deep Learning Models \& AI Frameworks}
\label{tab:tools_ch5}
\begin{tabularx}{\textwidth}{l l X}
\toprule
\textbf{Model / Tool} & \textbf{Framework} & \textbf{Official URL / Repository} \\
\midrule
\textbf{Python} & 3.10+ / 3.12 & \url{https://www.python.org} \\
\textbf{FastAPI} & ASGI Web Core & \url{https://fastapi.tiangolo.com} \\
\textbf{Uvicorn} & ASGI Server & \url{https://www.uvicorn.org} \\
\textbf{$\text{U}^2$-Net} & Background Removal & \url{https://github.com/xuebinqin/U-2-Net} \\
\textbf{rembg} & Segmentation Library & \url{https://github.com/danielgatis/rembg} \\
\textbf{TripoSR} & Stability AI & \url{https://github.com/VAST-AI-Research/TripoSR} \\
\textbf{LCM Dreamshaper v7} & Hugging Face Diffusers & \url{https://huggingface.co/SimianLuo/LCM_Dreamshaper_v7} \\
\textbf{Fal.ai API} & Cloud SD3.5 Large 2D & \url{https://fal.ai} \\
\textbf{Replicate API} & Cloud Stable Diffusion 3 & \url{https://replicate.com} \\
\bottomrule
\end{tabularx}
\end{table}

\subsection{Sequential VRAM Orchestration Logic}
Listing~\ref{lst:vram_orchestration} shows how PyTorch CUDA memory is safely allocated and purged sequentially in \texttt{main.py} to prevent kernel crashes under 8GB VRAM.

\begin{lstlisting}[language=Python, caption={Sequential VRAM Memory Purging Protocol (main.py)}, label={lst:vram_orchestration}]
@app.post("/api/summon")
def summon_protocol(req: SummonRequest):
    try:
        # Phase 1: Brain (Text-to-2D Image Generation)
        image = generate_image(req.prompt)
        torch.cuda.empty_cache() # Purge VRAM after 2D Tensor creation
        
        # Phase 2: Body (Image-to-3D Mesh Extraction)
        mesh = generate_3d(image)
        torch.cuda.empty_cache() # Purge VRAM after TripoSR extraction
        
        # Export GLB file and stream binary back to client
        temp_glb_path = export_temp_glb(mesh)
        with open(temp_glb_path, "rb") as f:
            glb_bytes = f.read()
        return Response(content=glb_bytes, media_type="model/gltf-binary")
    except Exception as e:
        raise HTTPException(status_code=500, detail=str(e))
\end{lstlisting}


% ==============================================================================
% CHAPTER 6: DATABASE, STORAGE & SESSION MANAGEMENT
% ==============================================================================
\section{Database, Storage \& Session Management}

\subsection{Database Architecture}
Vectra utilizes a dual-database architecture supporting both high-concurrency production setups and zero-configuration local deployments. Table~\ref{tab:tools_ch6} details database tooling.

\begin{table}[H]
\centering
\caption{Database Technologies \& Drivers}
\label{tab:tools_ch6}
\begin{tabularx}{\textwidth}{l l X}
\toprule
\textbf{Database Engine} & \textbf{Mode / Role} & \textbf{Official Website URL} \\
\midrule
\textbf{MySQL 8.0} & Production RDBMS (\texttt{sql_vectra_parsabe_com}) & \url{https://www.mysql.com} \\
\textbf{MariaDB} & Alternative Production Engine & \url{https://mariadb.org} \\
\textbf{SQLite 3} & Local Fallback (\texttt{database.sqlite}) & \url{https://www.sqlite.org} \\
\textbf{Laravel Breeze} & Auth \& Session Scaffolding & \url{https://laravel.com/docs/11.x/starter-kits} \\
\bottomrule
\end{tabularx}
\end{table}


% ==============================================================================
% CHAPTER 7: DOCKER CONTAINER & HOST INFRASTRUCTURE ECOSYSTEM
% ==============================================================================
\section{Docker Container \& Host Infrastructure Ecosystem}

\subsection{Host Web Server \& Nginx Configuration}
The production host environment runs Nginx under aaPanel/BT-Panel on Linux x86\_64 with Let's Encrypt TLS 1.2/1.3 SSL certificates.

\begin{table}[H]
\centering
\caption{Server Infrastructure \& Reverse Proxy Tools}
\label{tab:tools_ch7_infra}
\begin{tabularx}{\textwidth}{l l X}
\toprule
\textbf{Infrastructure Component} & \textbf{Role} & \textbf{Official URL} \\
\midrule
\textbf{Nginx Web Server} & Reverse Proxy \& Static Web Host & \url{https://nginx.org} \\
\textbf{Let's Encrypt} & Free SSL/TLS Certificates & \url{https://letsencrypt.org} \\
\textbf{aaPanel / BT-Panel} & Host Server Management & \url{https://www.aapanel.com} \\
\bottomrule
\end{tabularx}
\end{table}

\subsection{Docker Container Microservices}
The host machine runs several active Docker containers supporting adjacent AI pipelines and CI/CD operations (verified via \texttt{docker ps -a}). Table~\ref{tab:docker_containers} details these services.

\begin{table}[H]
\centering
\caption{Host Docker Container Ecosystem}
\label{tab:docker_containers}
\begin{tabularx}{\textwidth}{l l l X}
\toprule
\textbf{Container Name} & \textbf{Image Name} & \textbf{Port Mapping} & \textbf{Official Technology URL} \\
\midrule
\texttt{blackwall-service} & \texttt{blackwall-ai} & \texttt{8000:8000} & \url{https://fastapi.tiangolo.com} \\
\texttt{ollama-ollama-1} & \texttt{ollama/ollama:latest} & \texttt{11434:11434} & \url{https://ollama.com} \\
\texttt{jenkins-jenkins-1} & \texttt{jenkins/jenkins:2.462.1} & \texttt{14808:8080} & \url{https://www.jenkins.io} \\
\texttt{redis-redis-1} & \texttt{redis:7.4.0} & Internal & \url{https://redis.io} \\
\bottomrule
\end{tabularx}
\end{table}


% ==============================================================================
% CHAPTER 8: COMPREHENSIVE FILE & DIRECTORY INVENTORY
% ==============================================================================
\section{Comprehensive File \& Directory Inventory}

\begin{codecard}{Project File Tree Topology}
\begin{verbatim}
/www/wwwroot/vectra.parsabe.com
├── .env                                 # Environment variables (DB, API keys, Ports)
├── README.md / CHANGELOG.md             # Project overview and release history
├── package.json / composer.json         # NPM and PHP package dependencies
├── vite.config.js / tailwind.config.js  # Frontend asset compilation configs
├── main.py                              # Local GPU Engine (FastAPI + PyTorch + TripoSR)
├── proxy_server.py                      # VPS Proxy Gateway (FastAPI + Fal.ai/Replicate)
├── start.sh / stop.sh / status.sh       # Service management daemon scripts
├── app/                                 # Laravel Backend (Controllers, Models)
├── database/                            # Migrations & SQLite database file
├── resources/
│   ├── css/app.css                      # Cyberpunk Tailwind styling tokens
│   ├── js/app.js                        # Three.js 3D Viewport & Noclip Controls
│   ├── js/presentation.js               # 16-Stage Scrollytelling Presentation
│   ├── js/splat-sort.worker.js          # Web Worker for 3D Gaussian Splat sorting
│   └── views/                           # Blade templates (main, welcome, presentation)
├── TripoSR/                             # Stability AI TripoSR neural submodule
└── Paper codes/                         # Master's Thesis LaTeX source files (.tex/.bib)
\end{verbatim}
\end{codecard}


% ==============================================================================
% CHAPTER 9: DEPLOYMENT, EXECUTION & OPERATIONAL MANUAL
% ==============================================================================
\section{Deployment, Execution \& Operational Manual}

\subsection{Service Daemon Commands}
To start, check status, or stop backend servers:
\begin{lstlisting}[language=bash, caption={Server Daemon Execution Commands}]
# 1. Start all AI backend servers in background
./start.sh

# 2. Check operational status and PID health
./status.sh

# 3. Stop background servers
./stop.sh

# 4. Tail live execution logs
tail -f engine.log
tail -f proxy.log
\end{lstlisting}

\subsection{Offline CLI 3D Mesh Generation}
Generate 3D GLB models directly from terminal without web server overhead:
\begin{lstlisting}[language=bash, caption={Offline CLI Text-to-3D Mesh Generation}]
/www/wwwroot/venv/bin/python main.py "cyberpunk helmet"
# Outputs: output/output_<timestamp>.png and output/output_<timestamp>.glb
\end{lstlisting}


% ==============================================================================
% BIBLIOGRAPHY SECTION
% ==============================================================================
\newpage
\bibliographystyle{IEEEtran}
\bibliography{\jobname}

\end{document}
